\documentclass{../latex/utmreport}
\newcommand{\ReportTitle}{Lab 3: The Command Line: Shell, Pipes and Text Filters}
% Shared title-page and PDF metadata for every laboratory report.
\newcommand{\StudentName}{Gafenco Victor}
\newcommand{\StudentGroup}{FAF-241}
\newcommand{\CourseName}{Operating Systems}
\newcommand{\InstructorName}{Patricia Reițman}
\newcommand{\InstructorTitle}{Assistant lecturer}
\newcommand{\ReportCity}{Chișinău}
\newcommand{\ReportYear}{2026}
\newcommand{\CourseRepository}{https://github.com/yoda-dinmd/os-labs}

\selectlanguage{english}
\utmlogo{../latex/assets/UTM_Logo_English.jpeg}
\institution{Technical University of Moldova}
\faculty{Faculty of Computers, Informatics, and Microelectronics}
\department{Department of Software Engineering}
\studyprogramme{Software Engineering Study Programme}
\documenttype{Laboratory report}
\student{\StudentName}{\StudentGroup}
\supervisor{\InstructorName}{\InstructorTitle}
\renewcommand{\UTMCoordinatorLabel}{Instructor}
\cityyear{\ReportCity}{\ReportYear}
\reportsubtitle{\CourseName}
\hypersetup{pdfauthor={\StudentName},pdfsubject={\CourseName},
  pdftitle={\ReportTitle}}
\reporttitle{\ReportTitle}
\addto\captionsenglish{\renewcommand{\figurename}{Figure}\renewcommand{\tablename}{Table}\renewcommand{\bibname}{References}}
\renewcommand{\figurename}{Figure}
\renewcommand{\tablename}{Table}
\renewcommand{\bibname}{References}
\usepackage{float}
\usepackage{placeins}
\setlength{\emergencystretch}{2em}
\newcommand{\code}[1]{\texttt{\detokenize{#1}}}
% Keep screenshots beside their observations, with bounded height and width.
\newcommand{\screenshot}[4][\textwidth]{%
  \begin{figure}[H]
    \centering
    \includegraphics[width=#1,height=0.62\textheight,keepaspectratio]{screenshots/#2}
    \caption{#3}\label{#4}
  \end{figure}}
\newcommand{\reportreferences}[2]{%
  \clearpage\phantomsection\addcontentsline{toc}{chapter}{References}
  \begin{thebibliography}{9}
    \bibitem{handout} Technical University of Moldova, FCIM / FAF.
    \textit{#1}. Operating Systems laboratory handout, academic year 2026--2027.
    Supplied course document: \code{#2}.
  \end{thebibliography}}

\begin{document}
\makeutmreporttitlepage
\tableofcontents
\clearpage
\begin{utmabstract}[Abstract]
This work examines the shell as an interpreter between a user and operating-system
services. Experiments on openSUSE Tumbleweed establish how Bash resolves commands,
expands arguments, connects streams and manages asynchronous jobs. A deterministic
500-record data set supports eight pipeline exercises and a reusable reporting
script. Independent aggregation checks establish 100 failures and an even distribution
of requests across four pages. The results also demonstrate why a maximum-count
query must preserve ties and why redirection order changes observable behaviour.
\utmkeywords{Bash, expansion, file descriptors, pipes, processes, text filters}
\end{utmabstract}
\utmintroduction
The purpose of the laboratory is to explain command-line behaviour through
predictions and measured results, then combine small text-processing tools into a
working program. The course handout provides five practical parts~\cite{handout}.
The Linux exercises were run in \code{~/os-lab3}; the report and implementation
are retained in the existing course repository.

Commands use the C locale and Bash without personal startup files. Predictions
were saved before the first expansion experiments. Figures show saved terminal
output; the corresponding plain text is retained with the laboratory files.

\chapter{Command lookup and interactive assistance}
\section{Completion and history}
Tab completed \code{/usr/sh} to the unique existing directory \code{/usr/share/}.
The next Tab did not display competing alternatives in this environment.
Completion of \code{/et} produced \code{/etc/}, and executing the completed
\code{cd} command changed the working directory accordingly.
Ctrl+R recovered \code{echo history_marker_lab3} after two intervening commands;
Enter executed the recovered line. Completion and reverse search are shell editing
facilities, rather than features of the commands being prepared.
\screenshot{1.1-tab.png}{Completion targets and the resulting directory}{fig:l3-tab}

\section{Built-ins, PATH and failure status}
Both \code{cd} and \code{echo} resolved to Bash built-ins. The \code{ls} executable
resolved to \code{/usr/bin/ls}. A separate \code{cd} process could change only its
own current directory; therefore the interactive shell must perform this operation
in its own process. An external implementation of \code{echo} can coexist with the
built-in without changing the lookup result observed here.
\screenshot{1.2-type.png}{Command classification and executable location}{fig:l3-type}
The clean-shell PATH check contains five entries: \code{~/.local/bin},
\code{/usr/local/bin}, \code{/usr/bin}, \code{/bin} and \code{/usr/sbin}.
The count depends on the configured shell environment.
\screenshot{1.3-path.png}{PATH search locations and entry count}{fig:l3-path}
An unknown executable name produces status 127. The shell reports that it could
not find a command; this differs from a program starting successfully and then
returning an error of its own.
\screenshot{1.4-missing.png}{Status returned for a command that does not exist}{fig:l3-status}

\section{Finding help}
The requested commands were \code{man ls}, \code{ls --help},
\code{man -k "list directory"} and \code{help cd}. The first manual lookup
exposed a missing \code{ls} manual page on this installation. The corrected setup
includes \code{coreutils-doc}, the separate package for these pages. A final
lookup using the local manual search path succeeded and the manual index found
\code{ls}, \code{dir} and \code{vdir}. The option
\code{-t} sorts entries by modification time, newest first. Inside a manual pager,
\code{/sort} searches, \code{n} repeats the search, and \code{q} exits.
Bash's \code{help} documents its built-ins; \code{man -k} depends on the installed
manual-page index.

\chapter{Expansion and quoting}
\section{Experimental files and predictions}
The test directory initially contained two C filenames, three ordinary text
filenames, and \code{alfa beta.txt}. Later, brace expansion created three log files.
\screenshot{2.1-files.png}{Files used for expansion experiments}{fig:l3-files}
\begingroup\small\setstretch{1}
\begin{longtable}{p{0.32\textwidth}p{0.29\textwidth}p{0.30\textwidth}}
\textbf{Command} & \textbf{Prior prediction} & \textbf{Observed output}\\\hline
\code{echo *.c} & a.c b.c & a.c b.c\\
\code{echo file?.txt} & file1.txt file2.txt & file1.txt file2.txt\\
\code{echo [a-b]*} & a.c alfa beta.txt b.c & Same names, same order\\
\code{echo *.xyz} & Literal *.xyz & Literal *.xyz\\
\code{ls alfa beta.txt} & Two missing paths & Errors for alfa and beta.txt\\
\code{ls "alfa beta.txt"} & One filename; possible display quotes & 'alfa beta.txt'\\
\code{echo '$HOME' "$HOME" \$HOME} & Literal, expanded, literal & \code{$HOME /home/victor $HOME}\\
\code{echo {1..5}; echo $((7 * 6))} & 1 2 3 4 5; then 42 & As predicted\\
\lstinline!echo "today is $(date +%A), I am $(whoami)"! & Runtime weekday and login & today is Monday, I am victor\\
\code{touch log{1..3}.txt && ls log*} & Three numbered log names & log1.txt log2.txt log3.txt\\
\end{longtable}
\endgroup
The recorded predictions agreed with the semantic results. Terminal column spacing
and GNU \code{ls}'s presentation quotes do not change the underlying filenames.
The unmatched pattern stayed literal because default Bash glob options were used.
\screenshot{2.2-globs.png}{Matched and unmatched pathname patterns}{fig:l3-globs}

\section{Responsibility for constructing arguments}
The shell expands \code{*.c} before it invokes \code{echo}; the receiving program
need not enumerate a directory. A glob result containing a space remains one
argument, even though the displayed output alone does not show that boundary.
In contrast, unquoted spaces in the command text split \code{alfa beta.txt} into
two arguments. Either double quotes or a backslash before the space fixes the call.
\screenshot{2.3-spaces.png}{Failure caused by unquoted whitespace and two fixes}{fig:l3-spaces}
Single quotes inhibit expansion. Double quotes preserve word boundaries while
allowing variable, arithmetic and command substitution. An escaped dollar sign is
literal. GNU \code{ls} may add quotes when displaying a spaced filename in a
terminal; these are presentation characters, not part of the stored name.
\screenshot{2.4-quotes.png}{Quoting, arithmetic, brace and command expansion}{fig:l3-quotes}

\chapter{Streams and process control}
\section{Separating output and diagnostics}
In \code{ls /etc /nope > out.txt 2> err.txt}, descriptor 1 carries the directory
listing into \code{out.txt}, while descriptor 2 carries the missing-path diagnostic
into \code{err.txt}. Redirecting both streams to one file does not guarantee that
an error will appear at its end: buffering and write order still determine placement.
\screenshot{3.1-streams.png}{Directory output and error output in separate files}{fig:l3-streams}

\section{The order of redirections}
Redirection is processed from left to right. With \code{> all.txt 2>&1}, stdout
first changes destination and stderr then duplicates that destination. With
\code{2>&1 > all.txt}, stderr retains the original terminal destination while
stdout is subsequently moved to the file. The experiment therefore makes the
second version's error visible on the terminal.
\screenshot{3.2-order.png}{Different outcomes from reversing redirection order}{fig:l3-order}

\section{Input, chaining and asynchronous execution}
An input redirection gives \code{wc} bytes on standard input without a pathname
argument, so \code{wc -l < access.log} prints only 500. Supplying the pathname as
an argument additionally prints its name. Truncating and appending writes to
\code{notes.md} preserve the two lines \code{first} and \code{second}.

A semicolon sequences commands irrespective of status. The \code{&&} operator
continues only after success; \code{||} continues only after failure. The failing
GNU \code{ls /nope} returns 2, which the following conditional message reports.
\screenshot{3.3-operators.png}{Input redirection, appending and conditional commands}{fig:l3-operators}
The background \code{sleep 100} appears as job 1 and is signalled by
\lstinline!kill %1!. An immediate job listing may still show Running; a later check
observes termination. This illustrates asynchronous delivery and collection of
process status rather than a failed kill command.
\screenshot{3.4-jobs.png}{Starting, inspecting and terminating a background job}{fig:l3-jobs}

\chapter{Answering questions with pipelines}
\section{Data model and consistency checks}
The prescribed generator emits 500 records with date, time, user, result and page
fields. Every fifth index fails, every fourth selects the same page, and users
cycle modulo seven. Queries filter the relevant fields rather than matching
unrelated substrings. Each of Q1--Q8 occupies one executable line in
\code{oneliners.sh}; comments identify the question.

\section{Q1: failed requests}
The measured failure count is 100, consistent with one failure per five records.
\screenshot{4.1-q1.png}{Counting failed requests}{fig:l3-q1}
\section{Q2: distinct pages}
Sorting the fifth field and removing duplicates yields four paths, from
\code{/page0} through \code{/page3}.
\screenshot{4.2-q2.png}{The four distinct page names}{fig:l3-q2}
\section{Q3: requests per page}
Each page receives 125 requests. Their sum is 500, independently matching the
number of generated records.
\screenshot{4.3-q3.png}{Equal request counts across the pages}{fig:l3-q3}
\section{Q4: users with the most failures}
Both user3 and user5 have 15 failed requests. Every other user has 14. The query
retains every count equal to the first, maximum count after descending sorting;
using only \code{head -1} would incorrectly discard a tie. The seven counts sum
to 100, independently agreeing with Q1.
\screenshot{4.4-q4.png}{Both users sharing the maximum failure count}{fig:l3-q4}
\section{Q5: user3}
Filtering the third field yields 72 requests for user3, of which 15 failed.
\screenshot{4.5-q5.png}{Total and failed requests for user3}{fig:l3-q5}
\section{Q6: the last failed records}
The final three matching rows yield \code{10:10 user0}, \code{10:15 user5}, and
\code{10:20 user3}. Here last refers to file order: the generated minutes cycle
and must not be mistaken for a monotonically increasing timestamp.
\screenshot{4.6-q6.png}{Time and user from the last three failed records}{fig:l3-q6}
\section{Q7: configured account shells}
The observed account file contains two entries with \code{/bin/bash} and 32 with
\code{/usr/sbin/nologin}. These are account configurations, including service
accounts, rather than a count of current interactive users.
\screenshot{4.7-q7.png}{Counts grouped by the account shell field}{fig:l3-q7}
\section{Q8: the extra long-listing line}
The prior prediction was a difference of one. The two pipelines returned 213 and
214: ordinary piped \code{ls} emits one name per line, whereas its long form adds
a \code{total} header. These exact counts are specific to this installation.
Embedded newlines in filenames, aliases or concurrent changes would require
additional care when generalising the result.
\screenshot{4.8-q8.png}{The additional header in a long directory listing}{fig:l3-q8}
\screenshot{4.9-oneliners.png}{Executing all eight one-liners in order}{fig:l3-all}
\section{Regular-expression bonus}
A \code{grep -E} expression restricts the time to \code{10:[0-5]0} and the result
field to FAIL. It prints 50 records. Each tenth generator index satisfies the
minute condition and is also divisible by five. The executable query is retained
separately as \code{bonus.sh}.

\chapter{A reusable report program}
\section{Implementation and measured behaviour}
The completed script checks the argument count and file readability, obtains the
total and failed counts, then iterates over user0 through user6. Field-aware
aggregation returns explicit zeroes for empty sets. Quoted path arguments preserve
filenames containing spaces.
\screenshot{5.1-source.png}{Implementation of the reporting script}{fig:l3-source}
On the shared data it reports 500 requests and 100 failures, followed by counts
14, 14, 14, 15, 14, 15 and 14 for the seven users.
\screenshot{5.2-report.png}{Successful report for the generated log}{fig:l3-report}
A nonexistent pathname produces an error on descriptor 2 and status 1. Missing
or extra arguments also fail. Additional checks exercised an empty input, a
zero-failure input and a pathname containing spaces.
\screenshot{5.3-error.png}{Missing-file and missing-argument diagnostics}{fig:l3-error}

\section{Interpreter, permissions and exit status}
The shebang \code{#!/bin/bash} identifies the interpreter for direct execution.
Without execute permission, \code{./report.sh} is rejected, normally with status
126; explicitly running \code{bash report.sh} instead requires read permission.
The \code{./} prefix selects a pathname in the current directory rather than
relying on a PATH entry. Diagnostic output on stderr preserves the normal data
stream, and status 1 permits a caller to react to failure.

The optional cron bonus was not executed; no scheduled-run evidence is claimed.
The notebook includes a full-path example for that optional extension. Unrelated
crontab entries must be retained when adding or removing a temporary lab schedule.

\utmconclusions
The experiments identify three responsibilities that are easy to confuse: Bash
constructs arguments, redirections construct descriptor connections, and programs
consume the resulting arguments and streams. Small utilities become useful data
analysis tools when their inputs and fields are well defined. The agreement of
independent totals validates the pipelines, while the tied failure counts show
why seemingly minor query choices affect correctness. The final script makes
these observations reusable through explicit validation and meaningful status.
\end{document}
